\begin{definition}[Arithmetic Function]
    An \textit{arithmetic function} (also known as \textit{number-theoretic function}) is a function $f: \N \to S$, where $S = \Z$ or $S$ is a subset of $\C$ that is closed under addition and multiplication. 
\end{definition}

\section{Important Arithmetic Functions}
Unless stated otherwise, let prime factorization of $n \in \N$ be given by \[n = p_1^{\alpha_1} p_2^{\alpha_2} \cdots p_r^{\alpha_r},\] (where $p_i$ are distinct primes and $\alpha_i$ are positive integers) for the rest of this chapter. Some common arithmetic functions are:
\begin{enumerate}
    \item $\tau(n)$:  The function $\tau: \N \to \Z$ given by \[\tau(n) = \#\{d \in \N: d|n\}.\] is an arithmetic function, which basically counts the number of divisors of a given integer $n$. For example, $\tau(12) = 6$ because the divisors of $12$ are $1, 2, 3, 4, 6, 12$.

    \begin{prop}
        $\tau(n) = \prod_{i=1}^r (\alpha_i + 1)$.
    \end{prop}
    \begin{proof}
        Let $d = p_1^{\beta_1} p_2^{\beta_2} \cdots p_r^{\beta_r}$ be a divisor of $n$. Then, $0 \leq \beta_i \leq \alpha_i$ for all $i$. The number of choices for $\beta_i$ is $\alpha_i + 1$. Thus, the total number of divisors is given by the product of the number of choices for each prime factor, which is \[\prod_{i=1}^r (\alpha_i + 1).\]
    \end{proof}

    \item $\sigma(n)$: The function $\sigma: \N \to \Z$ given by \[\sigma(n) = \sum_{d|n} d\] is an arithmetic function, which basically sums the divisors of a given integer $n$. For example, $\sigma(12) = 1 + 2 + 3 + 4 + 6 + 12 = 28$.
    
    \begin{prop}\label{prop:sigmamul}
        $\sigma(n) = \prod_{i=1}^r \frac{p_i^{\alpha_i + 1} - 1}{p_i - 1}$.
    \end{prop}
    \begin{proof} 
        In order to evaluate $\sigma(n)$, consider the product \[(1 + p_1 + p_1^2 + \cdots + p_1^{\alpha_1})(1 + p_2 + p_2^2 + \cdots + p_2^{\alpha_2}) \cdots (1 + p_r + p_r^2 + \cdots + p_r^{\alpha_r}).\] The expansion of this product will yield all possible terms of the form $p_1^{\beta_1} p_2^{\beta_2} \cdots p_r^{\beta_r}$, where $0 \leq \beta_i \leq \alpha_i$. Each such term corresponds to a divisor of $n$ and all the divisors of $n$ can be obtained by varying the $\beta_i$'s. Hence, 
        \begin{align*}
            \sigma(n) &= (1 + p_1 + p_1^2 + \cdots + p_1^{\alpha_1})(1 + p_2 + p_2^2 + \cdots + p_2^{\alpha_2}) \\ &\cdots (1 + p_r + p_r^2 + \cdots + p_r^{\alpha_r})\\
            &= \frac{p_1^{\alpha_1 + 1} - 1}{p_1 - 1} \cdot \frac{p_2^{\alpha_2 + 1} - 1}{p_2 - 1} \cdots \frac{p_r^{\alpha_r + 1} - 1}{p_r - 1}\\
            &= \prod_{i=1}^r \frac{p_i^{\alpha_i + 1} - 1}{p_i - 1}.
        \end{align*}
    \end{proof}

    \item \textbf{Euler's Totient Function:} The function $\phi: \N \to \Z$ given by \[\phi(n) = \#\{k \in \N: 1 \leq k < n, \gcd(k, n) = 1\}\] is an arithmetic function, which counts the number of integers less than $n$ that are coprime to $n$. For example, $\phi(12) = 4$ because the integers less than $12$ that are coprime to $12$ are $1, 5, 7, 11$.
    
    Recall Corollary~\ref{cor:phi_form} which states \[\phi(n) = n \prod_{i=1}^r\left(1-\frac{1}{p_i}\right).\]

    \item $\epsilon(n)$: The function $\epsilon: \N \to \Z$ given by \[\epsilon(n) =1\] is also an arithmetic function.
    
    \item $\omega(n)$: The function $\omega: \N \to \Z$ given by \[\omega(n) = \#\{p_i\}\] is an arithmetic function, which counts the number of distinct prime factors of $n$. For example, $\omega(12) = 2$ because the prime factors of $12$ are $2$ and $3$.
\end{enumerate}

\section{Properties Arithmetic Functions}
Let $f$ be an arithmetic function and $n \in \N$, then \[\sum_{d \mid n}f(d) = \sum_{\frac{n}{d} \mid n}f(d) = \sum_{d \mid n}f\left(\frac{n}{d}\right).\]

Given arithmetic functions $f$ and $g$, we can define addition and multiplication pointwise, that is, for all $n \in \N$:
\begin{align*}
    (f + g)(n) &= f(n) + g(n), \\
    (f \cdot g)(n) &= f(n) \cdot g(n).
\end{align*}
\begin{prop}\label{prop:addarthfunc}
    The set of all arithmetic functions forms an abelian group under pointwise addition.
\end{prop}
\begin{proof}
    Let $f, g, h$ be any arithmetic functions. Then
    \begin{enumerate}
        \item \textbf{Closure under addition:} For all $n \in \N$, $(f + g)(n) = f(n) + g(n)$ is an integer, hence $f + g$ is an arithmetic function.
        \item \textbf{Identity element:} The function $\mathbf{0}(n) = 0$ is the identity element for addition, since \[(f + \mathbf{0})(n) = f(n) + \mathbf{0}(0) = f(n) + 0 = f(n).\] 
        \item \textbf{Inverse element:} For all $f$, the function $-f$ is the inverse element for addition, since \[(f + (-f))(n) = f(n) + (-f(n)) = 0.\]
        \item \textbf{Associativity:} For all $n \in \N$,
        \[((f+g)+h)(n) = (f+g)(n) + h(n)= f(n) + g(n) + h(n) \]\[= f(n) + (g+h)(n) = (f+(g+h))(n).\]
        \item \textbf{Commutativity:} For all $n \in \N$, \[(f + g)(n) = f(n) + g(n) = g(n) + f(n) = (g + f)(n).\]
    \end{enumerate}
    Hence, the set of all arithmetic functions forms an abelian group under addition.
\end{proof}
\begin{prop}
    The set of all arithmetic functions forms a ring under pointwise addition and multiplication.
\end{prop}
\begin{proof}
    We have already shown that the set of all arithmetic functions forms an abelian group under pointwise addition.
    We need to show that the set of all arithmetic functions is closed under multiplication and satisfies other ring axioms.
    Now, 
    \begin{enumerate}
        \item \textbf{Closure under multiplication:} For all $n \in \N$, $(f \cdot g)(n) = f(n) \cdot g(n)$ is an integer, hence $f \cdot g$ is an arithmetic function.
        \item \textbf{Associativity:} For all $n \in \N$,
        \[((f \cdot g) \cdot h)(n) = (f \cdot g)(n) \cdot h(n) = f(n) \cdot g(n) \cdot h(n)\]\[= f(n) \cdot (g \cdot h)(n) = (f \cdot (g \cdot h))(n).\]
        \item \textbf{Distributive Law:} For all $n \in \N$,
        \[(f \cdot (g + h))(n) = f(n) \cdot (g(n) + h(n)) \]\[= f(n) \cdot g(n) + f(n) \cdot h(n) = (f \cdot g + f \cdot h)(n).\]
        Similarly, we can show that \[((f + g) \cdot h)(n) = (f \cdot h + g \cdot h)(n).\]    
    \end{enumerate}
    Hence, the set of all arithmetic functions forms a ring under pointwise addition and multiplication.
\end{proof}

\section{Dirichlet Product and Its Properties}
There is another way of defining multiplication on arithmetic functions called \textit{Dirichlet product}. 
\begin{definition}[Dirichlet Product]
    Given arithmetic functions $f$ and $g$, the \textit{Dirichlet product} $\ast$ of $f$ and $g$ is defined as \[(f \ast g)(n) = \sum_{d\mid n} f(d) g\left(\frac{n}{d}\right).\]
\end{definition}

\begin{prop}\label{prop:dirprod}
    Given arithmetic functions $f$ and $g$, $h$, the Dirichlet product satisfies the following properties:
    \begin{enumerate}
        \item \textbf{Commutativity:} $f \ast g = g \ast f$.
        \item \textbf{Associativity:} $(f \ast g) \ast h = f \ast (g \ast h)$.
        \item \textbf{Identity element:} The function $\iota$ defined by \[\iota(n) = \begin{cases}
            1 & \text{if } n = 1 \\
            0 & \text{if } n > 1
        \end{cases}\] is the identity element for Dirichlet product.
        \item \textbf{Distributive Law:} $f \ast (g + h) = f \ast g + f \ast h$ and $(g + h) \ast f = g \ast f + h \ast f$.
    \end{enumerate}
\end{prop}
\begin{proof}
    \begin{enumerate}
        \item \begin{align*}
            (f \ast g)(n) &= \sum_{d\mid n} f(d) g\left(\frac{n}{d}\right) \\
            &= \sum_{d\mid n} g\left(\frac{n}{d}\right) f(d) \\
            &= \sum_{d\mid n} g(d) f\left(\frac{n}{d}\right) \\
            &= (g \ast f)(n).
        \end{align*}
    
    \item Writing \( d' = \frac{n}{d} \), we have
    \[
    (f * g)(n) = \sum_{dd' = n} f(d)g(d').
\]
Therefore
\begin{align*}
    ((f * g) * h)(n)
    &= \sum_{dd' = n} (f * g)(d) h(d') \\
    &= \sum_{dd' = n} \left( \sum_{d_1 d_2 = d} f(d_1)g(d_2) \right) h(d') \\
    &= \sum_{d_1 d_2 d' = n} f(d_1)g(d_2)h(d').
\end{align*}
\noindent
On the other hand,
\begin{align*}
    (f * (g * h))(n)
    &= \sum_{d'd = n} f(d') (g * h)(d) \\
    &= \sum_{d'd = n} f(d') \left( \sum_{d_1 d_2 = d} g(d_1) h(d_2) \right) \\
    &= \sum_{d' d_1 d_2 = n} f(d') g(d_1) h(d_2).
\end{align*}
Hence, \[((f \ast g) \ast h)(n) = (f \ast (g \ast h))(n).\]

\item \[(\iota \ast f)(n) = \sum_{d\mid n} \iota(d) f\left(\frac{n}{d}\right) = \iota(1) = 1,\]
since $\iota(d) = 0$ for all $d > 1$.

\item \begin{align*}
    (f \ast (g + h))(n) &= \sum_{d\mid n} f(d) (g + h)\left(\frac{n}{d}\right) \\
    &= \sum_{d\mid n} f(d) \left(g\left(\frac{n}{d}\right) + h\left(\frac{n}{d}\right)\right) \\
    &= \sum_{d\mid n} f(d) g\left(\frac{n}{d}\right) + \sum_{d\mid n} f(d) h\left(\frac{n}{d}\right) \\
    &= (f \ast g)(n) + (f \ast h)(n).
\end{align*}
Similarly, we can show that \[(g + h) \ast f = g \ast f + h \ast f.\]
\end{enumerate}
\end{proof}

\begin{remark}
    From proposition~\ref{prop:addarthfunc} and \ref{prop:dirprod}, it follows that the set of all arithmetic functions forms a Commutative ring (with zero divisors) under pointwise addition and Dirichlet product.  
\end{remark}

In a ring, not all elements have multiplicative inverses. However, for this particular ring, the following theorem states when an element has a multiplicative inverse.

\begin{theorem}\label{thm:dirprodinv}
    If $f(1) \neq 0$, then there exists a unique arithmetic function $h$ such that $f \ast h = \iota$.
\end{theorem}
\begin{proof}
    We will construct $h(n)$ for all $n \in \N$ recursively.

    For $n = 1$, we want \[(f \ast h)(1) = \iota(1) \Leftrightarrow  f(1)h(1) = 1 \Leftrightarrow h(1) = \frac{1}{f(1)}.\]

    Let $h(m) = \iota(m) = 0$ has been defined for all $m < n$. Then $h(n)$ is determined from the relation \[0 = \iota(n) = (f\ast h)(n) = f(1)h(n) + \sum_{\substack{d\mid n \\ d < n}} f\left(\frac{n}{d}\right)h(d)\]

    \[\implies h(n) = -\frac{1}{f(1)} \sum_{\substack{d\mid n \\ d < n}} f\left(\frac{n}{d}\right)h(d).\]

    Also the derivation determines $h(n)$ uniquely, completing the proof.
\end{proof}
\begin{remark}
    Using proposition~\ref{prop:dirprod} and theorem~\ref{thm:dirprodinv}, we conclude the set of all arithmetic functions satisfying the condition $f(1) \neq 0$ forms an abelian group under Dirichlet product.
\end{remark}